% Part: incompleteness
% Chapter: theories-computability
% Section: q-is-ce

\documentclass[../../../include/open-logic-section]{subfiles}

\begin{document}

\olfileid[hi]{inc}{tcp}{qce}

\olsection{$\Th{Q}$, \printtoken{S}{c.e.}-पूर्ण है}

\begin{thm}
$\Th{Q}$, !!{c.e.} है, किंतु निर्णेय नहीं है। वास्तव में वह एक पूर्ण
!!{c.e.} समुच्चय है।
\end{thm}

\begin{proof}
यह देखना कठिन नहीं है कि $\Th{Q}$, !!{c.e.} है, क्योंकि वह उन वाक्यों
(के कूटों)~$y$ का समुच्चय है जिनके लिए~$\Th{Q}$ में $y$ का कोई प्रमाण~$x$ है:
\[
Q = \Setabs{y}{\lexists[x][\Prf[\Th{Q}](x,y)]}.
\]
किंतु हम जानते हैं कि $\Prf[\Th{Q}](x,y)$ संगणनीय (वास्तव में आदिम
पुनरावर्ती) है, और ऊपर के रूप में लिखा जा सकने वाला प्रत्येक समुच्चय
संगणनीयतः परिगणनीय है।

उसके पूर्ण संगणनीयतः परिगणनीय समुच्चय होने का कथन $K \leq_m Q$ के समतुल्य है, जहाँ
$K = \Setabs{x}{!A_x(x) \downarrow}$। अतः दिखाएँ कि $K$ को~$\Th{Q}$ में
घटाया जा सकता है। चूँकि क्लेनी (Kleene) का विधेय $T(e,x,s)$ आदिम पुनरावर्ती
है, वह~$\Th{Q}$ में, मान लें $!A_T$ द्वारा, निरूपणीय है। तब प्रत्येक $x$
के लिए हमारे पास
\begin{align*}
x \in K & \lif \lexists[s][T(x,x,s)] \\
& \lif \lexists[s][(\Th{Q} \Proves !A_T(\num x,\num x, \num s))]
  \\
& \lif \Th{Q} \Proves \lexists[s][!A_T(\num x, \num x, s)].
\end{align*}
विलोमतः यदि
$\Th{Q} \Proves \lexists[s][!A_T(\num x, \num x, s)]$, तो वास्तव में
किसी प्राकृतिक संख्या $n$ के लिए सूत्र $!A_T(\num x, \num x, \num n)$
सत्य होना चाहिए। अब यदि $T(x,x,n)$ असत्य होता, तो $!A_T$ द्वारा $T$ का
निरूपण होने से $\Th{Q}$,~$\lnot !A_T(\num x, \num x, \num n)$ को सिद्ध
करता। किंतु तब $\Th{Q}$ एक असत्य सूत्र सिद्ध करता, जो विरोधाभास है। अतः
$T(x,x,n)$ सत्य होना चाहिए, जिससे $!A_x(x) \downarrow$ निकलता है।

संक्षेप में, प्रत्येक $x$ के लिए $x$,~$K$ में तभी और केवल तभी है जब
$\Th{Q}$,~$\lexists[s][T(\num x,\num x,s)]$ को सिद्ध करता है। अतः वह
फलन $f$ जो $x$ को वाक्य $\lexists[s][T(\num x, \num x, s)]$ के (किसी
कूट) पर भेजता है,~$K$ का~$\Th{Q}$ में एक अपचयन है।
\end{proof}

\end{document}
